\documentclass[11pt]{article}
% ===========================================================================
%  Shared preamble for the AP Calculus AB exam-preparation worksheets.
%  Input right after \documentclass. Files beginning with "_" are not compiled
%  on their own. Navy and gold, to match the course's lesson pages.
% ===========================================================================
\usepackage[letterpaper,top=2.3cm,bottom=2.2cm,left=1.7cm,right=1.7cm,headheight=28pt,footskip=30pt]{geometry}
\usepackage{amsmath,amssymb}
\usepackage[T1]{fontenc}
\usepackage[utf8]{inputenc}
\usepackage{lmodern}
\usepackage{xcolor}
\usepackage[most]{tcolorbox}
\usepackage{enumitem}
\usepackage{fancyhdr}
\usepackage{array}
\usepackage{tikz}
\usepackage{pgfplots}
\pgfplotsset{compat=1.18}

\definecolor{navy}{HTML}{161E54}
\definecolor{ink}{HTML}{27357E}
\definecolor{gold}{HTML}{B98A00}
\definecolor{goldbg}{HTML}{FFFDF3}
\definecolor{inkbg}{HTML}{F6F7FD}
\definecolor{rule}{HTML}{C9CFF0}
\definecolor{okgreen}{HTML}{1B7A44}
\definecolor{okbg}{HTML}{F1F8F3}

\pagestyle{fancy}
\fancyhf{}
\fancyhead[L]{\footnotesize Name: \underline{\hspace{4cm}}}
\fancyhead[C]{\textbf{\color{navy}Integration Academy}\\[-2pt]{\scriptsize Calculus AB \textperiodcentered{} AP\textsuperscript{\textregistered} Exam Preparation}}
\fancyfoot[L]{\tiny\color{black!55}\parbox[t]{0.78\textwidth}{AP\textsuperscript{\textregistered} is a registered trademark of the College Board, which was not involved in the production of, and does not endorse, this material. \textcopyright{} Integration Academy.}}
\fancyfoot[R]{\scriptsize\thepage}
\renewcommand{\headrulewidth}{1.2pt}
\renewcommand{\headrule}{\hbox to\headwidth{\color{gold}\leaders\hrule height \headrulewidth\hfill}}
\renewcommand{\footrulewidth}{0pt}
\setlength{\parindent}{0pt}
\setlength{\parskip}{4pt}

% Title block: \wstitle{code}{title}{unit line}
\newcommand{\wstitle}[3]{\begin{center}{\small\bfseries\color{gold}WORKSHEET #1}\\[3pt]{\LARGE\bfseries\color{navy}#2}\\[3pt]{\small\color{black!60}#3}\end{center}\vspace{2pt}}
% Section heading
\newcommand{\sech}[1]{\par\vspace{10pt}{\large\bfseries\color{navy}#1}\par\nobreak\vspace{2pt}{\color{rule}\hrule height0.8pt}\vspace{7pt}}
% A level tag printed beside a group of questions
\newcommand{\level}[1]{\par\vspace{6pt}{\footnotesize\bfseries\color{gold}\MakeUppercase{#1}}\par\nobreak\vspace{2pt}}

\newtcolorbox{ideabox}{colframe=ink,colback=inkbg,coltitle=white,colbacktitle=ink,fonttitle=\bfseries\small,title={Key ideas},arc=2pt,boxrule=0.8pt,left=8pt,right=8pt}
\newtcolorbox{example}[1]{colframe=rule,colback=white,coltitle=navy,colbacktitle=inkbg,fonttitle=\bfseries\small,title={#1},arc=2pt,boxrule=0.8pt,borderline west={3pt}{0pt}{ink},left=10pt,right=8pt}
\newtcolorbox{frbox}[1]{colframe=gold,colback=goldbg,coltitle=white,colbacktitle=gold,fonttitle=\bfseries\small,title={#1},arc=2pt,boxrule=0.8pt,breakable,left=8pt,right=8pt}
\newtcolorbox{solbox}{colframe=okgreen,colback=okbg,coltitle=white,colbacktitle=okgreen,fonttitle=\bfseries\small,title={Worked solutions},arc=2pt,boxrule=0.8pt,breakable,left=8pt,right=8pt}
\newtcolorbox{scorebox}[1]{colframe=navy,colback=white,coltitle=white,colbacktitle=navy,fonttitle=\bfseries\small,title={#1},arc=2pt,boxrule=0.8pt,breakable,left=8pt,right=8pt}

\newcommand{\soln}{\par\smallskip\textbf{\color{ink}Solution.}\ }
\newcommand{\result}[1]{\par\smallskip\textbf{#1}}
% blank working space of a given height
\newcommand{\work}[1]{\par\vspace{#1}}
% a point value in the margin of a scoring guide
\newcommand{\pts}[1]{\hfill{\footnotesize\bfseries\color{gold}[#1]}}

% Coordinate grid for graph questions: \begin{apgraph}{xmin}{xmax}{ymin}{ymax} ... \end{apgraph}
\newenvironment{apgraph}[4]{\begin{center}\begin{tikzpicture}\begin{axis}[
  axis lines=middle, xmin=#1, xmax=#2, ymin=#3, ymax=#4,
  grid=both, grid style={black!12}, xtick distance=1, ytick distance=1,
  ticklabel style={font=\scriptsize}, xlabel={$x$}, ylabel={$y$},
  width=10cm, height=7.2cm, samples=80, clip=false,
  every axis plot/.append style={line width=1.1pt, color=ink}]}{\end{axis}\end{tikzpicture}\end{center}}
\newcommand{\closeddot}[2]{\addplot[only marks, mark=*, mark size=2.2pt, color=ink] coordinates {(#1,#2)};}
\newcommand{\opendot}[2]{\addplot[only marks, mark=*, mark size=2.2pt, color=ink, mark options={fill=white}] coordinates {(#1,#2)};}

% Multiple-choice options: four across, or stacked when the options are long.
\newcommand{\opts}[4]{\par\vspace{3pt}\begin{tabular}{@{}p{0.235\linewidth}p{0.235\linewidth}p{0.235\linewidth}p{0.235\linewidth}@{}}(A)\ #1 & (B)\ #2 & (C)\ #3 & (D)\ #4\end{tabular}\par\vspace{9pt}}
\newcommand{\optsv}[4]{\par\vspace{3pt}\begin{tabular}{@{}l@{\ \ }p{0.9\linewidth}@{}}(A) & #1\\(B) & #2\\(C) & #3\\(D) & #4\end{tabular}\par\vspace{9pt}}

\fancyhead[R]{\footnotesize\color{black!60}Worksheet 5.11}
\begin{document}

\wstitle{5.11}{Behaviour of Implicitly Defined Curves}{Unit 5: Analysing Functions}

\begin{ideabox}
\begin{itemize}[leftmargin=*,itemsep=2pt,topsep=2pt]
  \item Write $\dfrac{dy}{dx}$ as a single fraction. The tangent is \textbf{horizontal} where the numerator is zero (and the denominator is not).
  \item The tangent is \textbf{vertical} where the denominator is zero (and the numerator is not).
  \item A condition such as $y=2x$ must be substituted back into the \textbf{original equation} to find actual points. Some conditions give no points.
  \item At a point with a horizontal tangent, the sign of $\dfrac{d^2y}{dx^2}$ decides between a maximum and a minimum.
\end{itemize}
\end{ideabox}

\sech{Worked examples}
\begin{example}{Example 1 \textperiodcentered{} Slope at a point}
Find the slope of $x^2+y^2=10$ at $(1,3)$.\soln $\dfrac{dy}{dx}=-\dfrac xy$.\result{$-\dfrac13$.}
\end{example}

\begin{example}{Example 2 \textperiodcentered{} Horizontal and vertical tangents}
Find the points on the ellipse $4x^2+y^2=16$ where the tangent is horizontal, and where it is vertical.\soln $8x+2y\dfrac{dy}{dx}=0$, so $\dfrac{dy}{dx}=-\dfrac{4x}{y}$. Numerator zero: $x=0$, $y=\pm4$. Denominator zero: $y=0$, $x=\pm2$.\result{Horizontal at $(0,\pm4)$; vertical at $(\pm2,0)$.}
\end{example}

\begin{example}{Example 3 \textperiodcentered{} Maximum or minimum}
For the same ellipse, find $\dfrac{d^2y}{dx^2}$ at $(0,4)$ and interpret it.\soln Differentiate $8x+2y\dfrac{dy}{dx}=0$ again: $8+2\left(\dfrac{dy}{dx}\right)^2+2y\dfrac{d^2y}{dx^2}=0$. At $(0,4)$ with $\dfrac{dy}{dx}=0$: $8+8\dfrac{d^2y}{dx^2}=0$.\result{$\dfrac{d^2y}{dx^2}=-1<0$: the curve has a relative maximum at $(0,4)$.}
\end{example}

\newpage
\sech{Practice}
\emph{Questions 5 to 8 refer to the curve $y^3-3y=x^2$.}
\level{Core}
\begin{enumerate}[leftmargin=*,itemsep=2pt]
  \item Find $\dfrac{dy}{dx}$ for $x^2+y^2=10$.\work{2.3cm}
  \item Find the slope of that circle at the point $(3,-1)$.\work{2.3cm}
  \item At which points of the circle $x^2+y^2=10$ is the tangent horizontal?\work{2.3cm}
  \item At which points of that circle is the tangent vertical?\work{2.3cm}
\end{enumerate}

\level{Exam level}
\begin{enumerate}[leftmargin=*,itemsep=2pt,start=5]
  \item Find $\dfrac{dy}{dx}$ in terms of $x$ and $y$.\work{2.8cm}
  \item Find all points of the curve at which the tangent line is horizontal.\work{2.8cm}
  \item Find all points of the curve at which the tangent line is vertical.\work{2.8cm}
  \item Verify that $\left(3\sqrt2,3\right)$ lies on the curve, and find the slope there.\work{2.8cm}
\end{enumerate}

\begin{frbox}{Question 9 \textperiodcentered{} Free-response practice}
Consider the curve $xy+y^2=8$.
\begin{enumerate}[label=(\alph*),leftmargin=*,itemsep=2pt]
  \item Find $\dfrac{dy}{dx}$ in terms of $x$ and $y$.\work{2.6cm}
  \item Write an equation of the line tangent to the curve at the point $(2,2)$.\work{2.6cm}
  \item Show that the curve has no horizontal tangent and no vertical tangent.\work{2.6cm}
\end{enumerate}
\end{frbox}

\newpage
\begin{solbox}
\begin{enumerate}[leftmargin=*,itemsep=4pt]
  \item $\dfrac{dy}{dx}=-\dfrac xy$.
  \item $-\dfrac{3}{-1}=3$.
  \item Where $x=0$: $\left(0,\sqrt{10}\right)$ and $\left(0,-\sqrt{10}\right)$.
  \item Where $y=0$: $\left(\sqrt{10},0\right)$ and $\left(-\sqrt{10},0\right)$.
  \item $\left(3y^2-3\right)\dfrac{dy}{dx}=2x$, so $\dfrac{dy}{dx}=\dfrac{2x}{3y^2-3}$.
  \item $x=0$ gives $y^3-3y=0$, so $y=0$ or $y=\pm\sqrt3$: the points $(0,0)$, $\left(0,\sqrt3\right)$ and $\left(0,-\sqrt3\right)$.
  \item $y=\pm1$. $y=1$ gives $x^2=-2$ (no points); $y=-1$ gives $x^2=2$: the points $\left(\sqrt2,-1\right)$ and $\left(-\sqrt2,-1\right)$.
  \item $27-9=18=\left(3\sqrt2\right)^2$. The slope is $\dfrac{6\sqrt2}{24}=\dfrac{\sqrt2}{4}$.
  \item \begin{enumerate}[label=(\alph*),leftmargin=*,itemsep=3pt]
    \item $y+x\dfrac{dy}{dx}+2y\dfrac{dy}{dx}=0$, so $\dfrac{dy}{dx}=-\dfrac{y}{x+2y}$.
    \item The slope is $-\dfrac{2}{6}=-\dfrac13$: $y-2=-\dfrac13(x-2)$.
    \item Horizontal needs $y=0$, but then $0=8$: no such point. Vertical needs $x=-2y$, which gives $-2y^2+y^2=8$, that is $y^2=-8$: no such point.
  \end{enumerate}
\end{enumerate}
\end{solbox}

\end{document}
